%%%%%%%%%%%%%%%%%%%%%%%%%%%%%%%%%%%%%%%%%%%%%%%%%%%%%%%%%%%%%%%
%        Template informe de trabajo de título ING PUC        %
%                  Creditos Github: @imant1                   %
%                                                             %
%          Creado a partir del template de práctica 2         %
%                 de Github: @diegocostares                   %
%%%%%%%%%%%%%%%%%%%%%%%%%%%%%%%%%%%%%%%%%%%%%%%%%%%%%%%%%%%%%%%



%%%%%%%%%%%%%%%%%%%%%%%%%%%%%%%%%%%%%%%%%%%%%%%%%%%%%%%%%%%%%%%
% IMPORTANTE: Cambiar el compilador a XeLaTeX en las opciones %
%     Si se quiere hacer "doble enter" usar: \vspace{2ex}     %
%%%%%%%%%%%%%%%%%%%%%%%%%%%%%%%%%%%%%%%%%%%%%%%%%%%%%%%%%%%%%%%
\documentclass{style}
\usepackage{tocloft}
\usepackage{adjustbox}

% Set page number at the top right corner
\clearpairofpagestyles
\ohead*{\thepage}


% La plantilla oficial centra los títulos de los índices con formato de título de capítulo (negrita, 12 pt)
\renewcommand{\cfttoctitlefont}{\hfill\normalsize\bfseries}
\renewcommand{\cftlottitlefont}{\hfill\normalsize\bfseries}
\renewcommand{\cftloftitlefont}{\hfill\normalsize\bfseries}

% La guía rotula la columna de páginas de cada índice con "Pág.", alineado a la derecha
\newcommand{\indexpagelabel}{\hfill\hfill\par\normalfont\normalsize\vspace{1ex}\hfill Pág.\par}
\renewcommand{\cftaftertoctitle}{\indexpagelabel}
\renewcommand{\cftafterlottitle}{\indexpagelabel}
\renewcommand{\cftafterloftitle}{\indexpagelabel}

% La guía lista las entradas como "I. TÍTULO" en el índice general y "Tabla 1.1: ..." / "Figura 1.1: ..." en los demás
\renewcommand{\cftsecaftersnum}{.}
\renewcommand{\cfttabpresnum}{Tabla~}
\renewcommand{\cfttabaftersnum}{:}
\renewcommand{\cftfigpresnum}{Figura~}
\renewcommand{\cftfigaftersnum}{:}
% Ancho reservado al número: debe caber la etiqueta más ancha, la de los anexos ("Tabla A.1:" / "Figura A.1:")
\cftsetindents{section}{0em}{2.6em}
\cftsetindents{table}{0em}{6em}
\cftsetindents{figure}{0em}{6.8em}


\begin{document}
\graphicspath{ {./img/} }

%%%%%%%%%%%% Numeración paginas %%%%%%%%%%%


%%%%%%%%%%%%%%%%% PORTADAS %%%%%%%%%%%%%%%%%

% La numeración sigue contando desde la portada, pero no se imprime aquí
\thispagestyle{empty}
\coverheader{%
  PONTIFICIA UNIVERSIDAD CATÓLICA DE CHILE\\
  ESCUELA DE INGENIERÍA
}

% Los bloques se separan con glue elástica (\vfill) en vez de medidas fijas: el
% sobrante de la página se reparte por igual, de modo que el conjunto
% título/autor queda centrado entre el encabezado y el bloque de la carrera.
% Ojo: tiene que ser \vfill (orden `fill') y no \stretch/\vfil, porque \newpage
% agrega su propio \vfil y se llevaría una parte del sobrante.
\vfill

\begin{coverbody}
  % El \par va DENTRO del grupo de \LARGE/\large: si se cierra antes, el
  % párrafo se rompe con el interlineado de 12pt y las líneas se solapan.
  \begin{center}
    {\LARGE\textbf{(TITULO TRABAJO)}\par}

    \vspace{2cm}
    {\large\textbf{(NOMBRE COMPLETO)}\par}
  \end{center}
\end{coverbody}

\vfill

\begin{coverbody}
  Trabajo de Título para optar al título de Ingeniero Civil en ...

  \vspace{1cm}

  Profesor Guía:\\
  \textbf{(NOMBRE PROFE)}
\end{coverbody}

\vfill

% La guía pide solo el año: "Santiago de Chile, (año)"
\begin{coverbody}
  Santiago de Chile, [AÑO]
\end{coverbody}
% Deja la fecha un poco por encima del punto donde termina la línea
% vertical (a 2,5 cm del borde inferior), como en la plantilla oficial.
\vspace{0.5cm}
\newpage


% The counter keeps advancing even though the number is not printed here
\thispagestyle{empty}
\coverheader{%
  PONTIFICIA UNIVERSIDAD CATÓLICA DE CHILE\\
  ESCUELA DE INGENIERÍA\\
  Departamento de [COMPLETAR]
}

\vfill

\begin{coverbody}
  \begin{center}
    {\LARGE\textbf{(TITULO TRABAJO)}\par}

    \vspace{2cm}
    {\large\textbf{(NOMBRE COMPLETO)}\par}
  \end{center}
\end{coverbody}

\vfill

\begin{coverbody}
  Trabajo de Título presentado a la Comisión integrada por los profesores:

  \textbf{(PROFESOR GUÍA)}\\
  \textbf{(PROFESOR REPRESENTANTE DE PREGRADO)}\\
  \textbf{(PROFESOR INVITADO)}\\
  \textbf{(PROFESOR INVITADO)}

  \vspace{1cm}

  Para completar las exigencias del título de
  Ingeniero Civil en [COMPLETAR]
\end{coverbody}

\vfill

% La guía pide solo el año: "Santiago de Chile, (año)"
\begin{coverbody}
  Santiago de Chile, [AÑO]
\end{coverbody}
\vspace{0.5cm}
\newpage

%%%%%%%%%%%%%%%%% DEDICATORIA %%%%%%%%%%%%%%%%%
% Dedication Page
\vspace*{0.7\textheight}
\hfill
\tocsection{DEDICATORIA}
\begin{minipage}{0.4\textwidth}
    (A mis Padres, hermanos y amigos,\\
    que me apoyaron mucho.)\\
    \textit{(Dedicatoria)}
\end{minipage}
\newpage

%%%%%%%%%%%%%%%%% AGRADECIMIENTOS %%%%%%%%%%%%%%%%%
% Acknowledgments Section
\begin{center}
    \textbf{AGRADECIMIENTOS}
\end{center}
\tocsection{AGRADECIMIENTOS}
\vspace{1cm}
(Nota redactada sobriamente en la cual se agradece a quienes han colaborado en la elaboración del trabajo.) \textit{(Normal)}

\newpage

%%%%%%%%%%%%%%%%% INDICES %%%%%%%%%%%%%%%%%
% En la guía el índice general no se lista a sí mismo: parte en DEDICATORIA y sigue en ÍNDICE DE TABLAS
\renewcommand{\contentsname}{ÍNDICE GENERAL} % Center the title
\tableofcontents \newpage

% La guía exige el índice de tablas, igual que el de figuras
\renewcommand{\listtablename}{ÍNDICE DE TABLAS}
\tocsection{ÍNDICE DE TABLAS}
\listoftables \newpage

\renewcommand{\listfigurename}{ÍNDICE DE FIGURAS}
\tocsection{ÍNDICE DE FIGURAS}
\listoffigures \newpage



%%%%%%%%%%%%%%%%% RESUMEN %%%%%%%%%%%%%%%%%
%%%%% ABSTRACT %%%%


\begin{center}
    \textbf{RESUMEN}
\end{center}
\tocsection{RESUMEN}
\label{sec:RESUMEN}
\vspace{1cm}


\newpage

%%%%%%%%%%%%%%%%% ABSTRACT %%%%%%%%%%%%%%%%%
%%%%% ABSTRACT %%%%

\selectlanguage{British}

\begin{center}
    \textbf{ABSTRACT}
\end{center}
\tocsection{ABSTRACT}
\vspace{1cm}
(Abstract en inglés que no debe exceder una página.) (Normal)
\selectlanguage{Spanish}
\newpage

%%%%%%%%%%%%%%%%% CONTENIDO %%%%%%%%%%%%%%%%%

%%%% Estas secciones son a modo de ejemplo %%%%

\section{INTRODUCCIÓN}

IMPORTANTE: estas secciones son a modo de ejemplo, no lo tomes como que debe ser así.

\subsection{Descripción de la empresa}

\begin{table}[H]
    \centering
    \caption{Comparación de técnicas basadas en modelos}
    \label{tab:ejemplo}
    \begin{tabular}{lcc}
        \hline
        Concepto              & Modelación analítica & Modelación por simulación \\
        \hline
        Tiempo de aprendizaje & 4                    & 1                         \\
        Costo de construcción & 1                    & 5                         \\
        \hline
    \end{tabular}
\end{table}

En la Tabla \ref{tab:ejemplo} se muestra un ejemplo de tabla con su rótulo y su entrada en el índice de tablas.

\subsection{Rol del estudiante}

\subsection{Objetivos del proyecto}



\subsection{Competencias de egreso}

\newpage
\section{TRABAJO REALIZADO POR EL ESTUDIANTE}
\subsection{Metodología de trabajo} 
\subsection{Desarrollo técnico del proyecto}



\subsection{Resultados obtenidos}


\newpage
\section{COMPETENCIAS DEMOSTRADAS}

\newpage
\section{CONCLUSIONES}


%%%%%%%%%%%%%%%%% BIBLIOGRAFIA %%%%%%%%%%%%%%%%%
\newpage
% Hay 2 formas de agregar bibliografia:
% 1. Agregar una bibliografia en un archivo .bib (Es super automatico)
%    En el preambulo: \usepackage[style=apa]{biblatex} y \addbibresource{mybib.bib}
%    La guia titula esta seccion "BIBLIOGRAFIA", sin numeral de capitulo.
% \printbibliography[title={BIBLIOGRAFIA}] % Requiere compilar con biber

% 2. Agregar una bibliografia en un archivo .tex
% (Es manual, pero comodo para los que no conoce el bibtex)
\indent
% La guia lista la bibliografia sin numeral: no entra en la secuencia de
% capitulos (I., II., ...). Va como seccion sin numerar y se agrega al indice a
% mano, igual que RESUMEN, ABSTRACT o GLOSARIO.
\section*{BIBLIOGRAFIA}
\tocsection{BIBLIOGRAFIA}
\leftskip 0.3in
\parindent -0.3in

%%%%%%%%%%%%%%% ESCRIBIR DEBAJO EN FORMATO APA %%%%%%%%%%%%%%%
% las lineas de arriba permiten que tenga una sangría francesa.
 \newpage

%%%%%%%%%%%%%%%%% ANEXOS %%%%%%%%%%%%%%%%%
%%%%% GLOSARIO %%%%


\begin{center}
    \textbf{GLOSARIO}
\end{center}
\tocsection{GLOSARIO} 
\vspace{1cm}

(opcional)


\newpage

\begin{center}
\vspace*{\fill}
    
% La guia titula la portadilla "A N E X O S". El espaciado lo pone
% \spacedletters (soul) y no espacios literales, para que la palabra siga
% siendo una sola. En el indice va sin espaciar: la palabra separada se lee
% mal en una lista de titulos.
\textbf{\spacedletters{ANEXOS}}
\tocsection{ANEXOS}
\vspace*{\fill}
\end{center}
\newpage



%%%% para agregar sectiones o subsecciones al anexo ocupar 
%%%% \addannexsec{--title name--}
%%%% \addannexsubsec{--title name--}
\startannex{A}{Panel de reportería de finanzas}



\end{document}
